% !TEX program = tectonic
% FICHE DONNÉES HOUSE — A4 portrait, 2 colonnes, 9 étapes en 2 actes, 3 pages pleines.
% Source reconstituée à l'identique du PDF livré (le .tex d'origine a disparu avec le worktree v1-house,
% supprimé au commit fdafda12 : seul le PDF avait été rapatrié).
% RÈGLES DE FORME (ne pas dévier) :
%   · corps en \footnotesize partout (texte, puces, tableaux) ; \scriptsize gris pour les seules
%     sources ; un unique chiffre-phare en \Large par étape, jamais d'autre variation de taille.
%   · une étape = une boîte insécable au gabarit constant : titre + ruban de position, une phrase
%     qui répond, le chiffre-phare et son unité, deux à trois puces, un visuel, la source.
%   · aucune information nécessaire à la compréhension en \src ; phrases courtes, une idée par phrase.
%   · zéro vocabulaire de performance : on décrit la donnée et on s'arrête avant toute stratégie.
% Chaque chiffre sort d'une cellule exécutée de Portefeuilles_House_Complet.ipynb.
\documentclass[10pt,a4paper]{article}
\usepackage[margin=1.05cm]{geometry}
\usepackage[T1]{fontenc}
\usepackage[utf8]{inputenc}
\usepackage{lmodern}
\usepackage[french]{babel}
\usepackage{amsmath,amssymb}
\usepackage{booktabs}
\usepackage{multicol}
\usepackage{xcolor}
\usepackage{graphicx}
\usepackage{enumitem}
\usepackage{tikz}
\usepackage[most]{tcolorbox}
\usetikzlibrary{positioning}

\definecolor{bleu}{RGB}{20,77,144}
\definecolor{vert}{RGB}{20,110,60}
\definecolor{orange}{RGB}{196,123,0}
\definecolor{rouge}{RGB}{160,30,30}
\definecolor{grisf}{RGB}{110,110,110}
\definecolor{grisc}{RGB}{205,205,205}

\setlength{\columnsep}{5.5mm}
\setlength{\parindent}{0pt}
\setlength{\parskip}{0pt}
\pagestyle{empty}
\setlist[itemize]{leftmargin=3.4mm,topsep=0.9mm,itemsep=0.85mm,parsep=0pt,label=\textbullet}

% ── le ruban de position : 9 pastilles, la courante pleine ────────────────────
\newcommand{\ruban}[1]{%
  \begin{tikzpicture}[baseline=-0.6ex, x=2.35mm, y=2.35mm]
    \foreach \k in {1,...,9}{%
      \ifnum\k=#1\relax
        \fill[white] (\k,0) circle (0.9);
      \else
        \fill[white, opacity=0.42] (\k,0) circle (0.6);
      \fi}
  \end{tikzpicture}}

% ── LE gabarit d'étape ───────────────────────────────────────────────────────
\newtcolorbox{etapebox}[3]{%
  enhanced, breakable=false, colback=#2!3, colframe=#2!55, boxrule=0.5pt,
  arc=2pt, left=2.2mm, right=2.2mm, top=1.8mm, bottom=1.8mm,
  toptitle=0.7mm, bottomtitle=0.7mm, boxsep=0pt,
  title={\begin{minipage}[c]{0.63\linewidth}\raggedright\footnotesize\bfseries
           \textcolor{white}{#3}\end{minipage}\hfill
         \begin{minipage}[c]{0.30\linewidth}\raggedleft\ruban{#1}\end{minipage}},
  colbacktitle=#2, coltitle=white, fonttitle=\footnotesize\bfseries,
  before skip=1.9mm, after skip=0mm}

% ── les briques internes, invariables ────────────────────────────────────────
\newcommand{\reponse}[1]{{\footnotesize #1\par}\vspace{1.1mm}}
\newcommand{\phare}[3]{%          #1 = chiffre  #2 = unité  #3 = couleur
  \begin{center}\vspace{-0.6mm}
  {\Large\bfseries\color{#3} #1}\\[0.4mm]{\footnotesize #2}
  \end{center}\vspace{0.6mm}}
\newcommand{\pts}[1]{{\footnotesize\begin{itemize}#1\end{itemize}}\vspace{0.9mm}}
\newcommand{\vis}[1]{\vspace{0.8mm}\begin{center}\includegraphics[width=0.94\linewidth]{figures/#1}\end{center}\vspace{0.4mm}}
\newcommand{\src}[1]{{\scriptsize\color{grisf} #1\par}}
\newcommand{\acte}[2]{%
  \vspace{1.1mm}{\footnotesize\bfseries\color{#2} #1}\vspace{-1.3mm}
  \par{\color{#2}\rule{\linewidth}{1pt}}\vspace{0.6mm}}
% ── la grille : deux boîtes côte à côte, alignées en haut ; pagination explicite ──
\newlength{\colw}\setlength{\colw}{0.4865\textwidth}
\newcommand{\rangee}[2]{%
  \noindent\begin{minipage}[t]{\colw}\vspace{0pt}#1\end{minipage}%
  \hfill\begin{minipage}[t]{\colw}\vspace{0pt}#2\end{minipage}\par}

\begin{document}

% ═══════════════ CHAPEAU PLEINE LARGEUR ═══════════════
\begin{center}
{\large\bfseries La donnée House, étape par étape}\\[0.5mm]
{\scriptsize ce que les élus déposent, ce qu'on en reconstruit, comment on le vérifie
$\cdot$ House 2014-2026 $\cdot$ notebook \texttt{Portefeuilles\_House\_Complet}}
\end{center}
\vspace{1.4mm}

\begin{tcolorbox}[enhanced, colback=white, colframe=grisc, boxrule=0.5pt, arc=2pt,
                  left=3mm, right=3mm, top=1.5mm, bottom=1.5mm]
{\footnotesize\bfseries\color{bleu} ACTE I --- la vie d'un mandat}\hfill
{\footnotesize\bfseries\color{vert} ACTE II --- ce qu'on en fait}\\[0.9mm]
\begin{center}
\begin{tikzpicture}[x=0.98cm, y=0.9cm, font=\scriptsize]
\foreach \k/\lab/\c [count=\i from 0] in {%
  1/il arrive/bleu, 2/il trade/bleu, 3/il déclare/bleu, 4/il part/bleu,
  5/on rassemble/vert, 6/on trie/vert, 7/on valorise/vert, 8/on reconstruit/vert, 9/on vérifie/vert}{%
  \ifnum\k<9 \draw[\c!45, line width=0.9pt] ({\i*2.1+0.28}, 0) -- ({\i*2.1+1.82}, 0);\fi
  \fill[\c] ({\i*2.1}, 0) circle (0.23);
  \node[text=white, font=\scriptsize\bfseries] at ({\i*2.1}, 0) {\k};
  \node[text=\c, font=\scriptsize, anchor=north] at ({\i*2.1}, -0.30) {\lab};}
\end{tikzpicture}
\end{center}
\vspace{0.2mm}
{\footnotesize\textbf{Les mots, une fois pour toutes.} Chaque nombre de cette fiche dit laquelle de ces
quatre choses il compte.}\\[1.4mm]
\begin{center}
\begin{tikzpicture}[x=1cm, y=1cm, font=\scriptsize,
  bo/.style 2 args={rectangle, rounded corners=1.8pt, draw=#1, fill=#1!8, semithick,
                    inner sep=2.8pt, align=center, text width=#2, font=\scriptsize}]
\node[bo={grisf}{1.9cm}]  (m) at (0, 0)    {un \textbf{membre}};
\node[bo={bleu}{2.6cm}]   (d) at (3.1, 0)  {dépose des \textbf{documents}};
\node[bo={vert}{2.6cm}]   (l) at (6.5, 0)  {qui contiennent des \textbf{lignes}};
\node[bo={orange}{4.6cm}] (p) at (11.4, 0.62) {une \textbf{position} détenue, dans une photo};
\node[bo={orange}{4.6cm}] (t) at (11.4, -0.62) {une \textbf{transaction} datée, dans un flux};
\draw[->, grisf, semithick] (m) -- (d);
\draw[->, grisf, semithick] (d) -- (l);
\draw[->, grisf, semithick] (l.east) -- ++(0.5, 0) |- (p.west);
\draw[->, grisf, semithick] (l.east) -- ++(0.5, 0) |- (t.west);
\node[font=\scriptsize\itshape, text=grisf] at (8.9, 0.98) {chaque ligne est};
\end{tikzpicture}
\end{center}
\vspace{1mm}
{\footnotesize La \textbf{reconstruction} est le calcul qui rejoue ces lignes pour retrouver, jour après
jour, ce que chaque membre détenait.}
\end{tcolorbox}

\vspace{1.6mm}
% ═══════════════ CORPS : bandeau population, puis les 9 étapes en rangées ═══════════════
\begin{tcolorbox}[enhanced, breakable=false, colback=grisc!12, colframe=grisc, boxrule=0.5pt,
                  arc=2pt, left=2.2mm, right=2.2mm, top=1.8mm, bottom=1.8mm,
                  toptitle=0.7mm, bottomtitle=0.7mm, boxsep=0pt, before skip=1mm, after skip=0mm,
                  colbacktitle=grisf, coltitle=white, fonttitle=\footnotesize\bfseries,
                  title={LA POPULATION --- 900 membres, 4 trajectoires}]
\reponse{Quatre obligations rythment un mandat. Ce qu'on attend d'un membre dépend donc entièrement de
sa trajectoire.}
\noindent\begin{minipage}[t]{0.485\linewidth}\vspace{0pt}\vis{fig_10_obligations.png}\end{minipage}%
\hfill\begin{minipage}[t]{0.485\linewidth}\vspace{0pt}\vis{fig_10_quadrant.png}\end{minipage}
\pts{
\item Un \textbf{arrivé} doit une photo d'entrée. Un \textbf{partant} doit une photo de sortie. Les deux
à la fois pour les 149 mandats complets, qui sont les seuls observables de bout en bout.
\item Ceux qui étaient déjà en poste en 2014 sont entrés avant notre fenêtre : on n'attend d'eux aucune
photo d'entrée.
\item Cette répartition, construite sur nos dépôts, concorde à 98,9\,\% avec le référentiel officiel des
mandats. Les 10 cas de frontière sont listés dans le notebook.
}
\src{source : cellules \og mandats officiels \fg{} et \og validation croisée \fg{}.}
\end{tcolorbox}

\vfill
\rangee{\begin{etapebox}{1}{bleu}{1 $\cdot$ IL ARRIVE --- la photo d'entrée}
\reponse{En prenant ses fonctions, l'élu déclare tout son patrimoine. C'est notre point de départ : sans
cette photo, on ne sait pas ce qu'il possédait avant de commencer à trader.}
\phare{200 sur 312}{arrivés encore en poste ayant déposé leur photo d'entrée}{bleu}
\vis{fig_10_etape_arrivee.png}
\pts{
\item Chez les mandats complets (arrivés \emph{et} repartis dans la fenêtre), c'est 80 sur 149.
\item Les 181 photos manquantes sont autant de membres dont le point de départ nous échappe. Parmi celles
déposées, toutes ne seront pas utilisables : l'étape 8 le montre.
\item La photo n'arrive pas le jour de l'entrée : il s'écoule 210 jours en médiane entre la prise de
fonction et le dépôt.
}
\src{source : cellules \og mandats officiels \fg{} et \og tableau central \fg{} du notebook.}
\end{etapebox}}{\begin{etapebox}{2}{bleu}{2 $\cdot$ IL TRADE --- la déclaration au fil de l'eau}
\reponse{Chaque achat ou vente doit être déclaré dans les jours qui suivent. C'est le document le plus
rapide : 27 jours en médiane entre la transaction et sa publication, donc connu presque tout de suite.}
\phare{391 sur 900}{membres qui tradent des actions cotées}{bleu}
\vis{fig_10_trois_zones.png}
\pts{
\item 305 membres sont vus par leurs déclarations au fil de l'eau, dont 199 aussi par leurs rapports annuels.
\item 86 membres ne sont vus \textbf{que} par leurs rapports annuels : ils n'ont jamais déposé de
déclaration exploitable. Un cinquième des traders est donc invisible au document rapide.
\item Les 509 autres ne tradent pas d'actions cotées. Ils déclarent pourtant : 444 d'entre eux déposent
leurs rapports annuels. Ne pas trader n'est pas ne pas déclarer.
\item 315 membres ont déclaré au moins une transaction. Pour dix d'entre eux, aucun prix n'est
retrouvable : ils sortent du décompte, d'où les 305.
}
\src{source : cellules \og qui trade par trajectoire \fg{} et \og trois zones \fg{}.}
\end{etapebox}}
\vfill
\newpage

\rangee{\begin{etapebox}{3}{bleu}{3 $\cdot$ IL DÉCLARE --- le rapport annuel}
\reponse{Chaque année de mandat donne lieu à un rapport : le patrimoine au 31 décembre, plus la liste des
transactions de l'année. C'est la pièce qui nous permettra de vérifier notre travail.}
\phare{85 \%}{des années de mandat couvertes par un rapport (médiane par membre)}{bleu}
\vis{fig_10_etape_annuel.png}
\pts{
\item Sur toute la fenêtre, 5\,702 rapports sont attendus : 212 membres couvrent toutes leurs années,
79 n'en couvrent aucune.
\item La photo de sortie tient lieu de dernier rapport annuel : elle est comptée comme telle ici, ainsi
que les amendements, ces corrections déposées après coup.
\item Les transactions listées dans ces rapports sortent tard : 373 jours en médiane après avoir eu lieu.
On les découvre donc plus d'un an après coup.
}
\src{source : cellule \og couverture des rapports annuels \fg{}. Années attendues calculées depuis les
mandats officiels : un rapport est dû pour chaque année où le membre est en poste au 31 décembre.}
\end{etapebox}}{\begin{etapebox}{4}{bleu}{4 $\cdot$ IL PART --- la photo de sortie}
\reponse{En quittant ses fonctions, l'élu dépose une dernière photo de son patrimoine. C'est le point
d'arrivée qui permet de boucler une trajectoire complète.}
\phare{253 sur 314}{partants déjà en poste en 2014 ayant déposé leur photo de sortie}{bleu}
\vis{fig_10_etape_sortie.png}
\pts{
\item Chez les mandats complets, c'est 114 sur 149.
\item Ces 149 mandats complets sont les seuls à avoir à la fois un début et une fin dans nos données.
\item Le départ est suivi vite : 34 jours en médiane entre la fin du mandat et le dépôt.
}
\src{source : cellules \og tableau central \fg{} et \og délais administratifs \fg{}.}
\end{etapebox}}
\begin{tcolorbox}[enhanced, breakable=false, colback=bleu!4, colframe=bleu!45, boxrule=0.5pt,
                  arc=2pt, left=2.4mm, right=2.4mm, top=1.4mm, bottom=1.4mm,
                  before skip=1.9mm, after skip=0mm]
{\footnotesize\textbf{Fin de l'acte I.} On sait qui devait déposer quoi, et ce qu'on a reçu.
\textbf{Une même transaction, deux horloges} : déclarée au fil de l'eau, elle est publique en 27 jours ;
reprise dans le rapport annuel, elle ne l'est qu'après 373 jours. Reste à ouvrir ces documents.}
\end{tcolorbox}

\acte{ACTE II --- ce qu'on en fait : de la ligne déclarée au portefeuille vérifié}{vert}
\rangee{\begin{etapebox}{5}{vert}{5 $\cdot$ ON RASSEMBLE --- l'inventaire des documents}
\reponse{Premier travail : réunir tous les documents déposés et compter ce qu'ils contiennent, sans encore
rien en écarter.}
\phare{420\,916}{lignes de patrimoine, dans 6\,135 photos}{vert}
\vis{fig_10_inventaire.png}
\pts{
\item Les photos se répartissent en 4\,447 rapports annuels, 934 photos de candidat, 378 de sortie et
376 d'entrée --- ces dernières se ramenant à 284 photos retenues, une par membre.
\item Deux contrôles de bouclage : 18 membres du registre n'ont laissé aucun document, et 12 déposants sont
des candidats jamais élus, donc hors registre.
}
\src{source : cellules \og inventaire documentaire \fg{} et \og bouclage \fg{}. Les dépôts de déclarations
rapides sont comptés en paires (membre, date de publication) : leur identifiant de document n'existe pas
dans la table.}
\end{etapebox}}{\begin{etapebox}{6}{vert}{6 $\cdot$ ON TRIE --- chaque ligne reçoit une classe}
\reponse{Toutes les lignes déclarées ne sont pas des actions. Avant de chercher le moindre prix, chaque
ligne est classée : action cotée, fonds mutuel, monétaire, ou faux ticker.}
\phare{156\,516 sur 420\,916}{lignes qui sont vraiment des actions ou des ETF}{vert}
\vis{fig_10_paysage.png}
\pts{
\item La majorité de ce qui est déclaré n'a pas de ticker : immobilier, sociétés non cotées, comptes.
Cela pèse 68,3 Md\$ contre 10,6 Md\$ d'actions et d'ETF.
\item \textbf{La règle qu'on s'impose} : le moteur ne voit que des actions et des ETF. Le reste est compté,
jamais valorisé comme une action --- y compris les faux tickers, ces codes de formulaire (EIF, SP,
DC) qu'une lecture automatique confond facilement avec des symboles boursiers.
\item Le piège que cette règle ferme : côté transactions (étape 7), 8\,439 des 10\,971 lignes écartées
par leur classe ont pourtant un prix disponible, parce qu'un fonds a une valeur liquidative. Un filtre
\og a-t-il un prix ? \fg{} les aurait laissées entrer.
}
\src{source : cellules \og classification \fg{} et \og paysage déclaré \fg{}. Valeurs = milieu de fourchette déclarée.}
\end{etapebox}}
\vfill
\newpage

\vspace*{\fill}
\rangee{\begin{etapebox}{7}{vert}{7 $\cdot$ ON VALORISE --- les transactions retenues}
\reponse{Une transaction n'entre dans le moteur que si l'on connaît le prix de son titre le jour dit. Voici
ce qui reste, et pourquoi le reste tombe.}
\phare{121\,300}{transactions datées et valorisées}{vert}
\begin{center}\footnotesize
\begin{tabular}{@{}lrr@{}}
\toprule
 & \color{bleu}au fil de l'eau & \color{orange}rapports annuels \\
\midrule
lignes déclarées & 122\,814 & 26\,115 \\
$-$ sens, dates, montants & $-5$ & $-235$ \\
$-$ classe non négociable & --- & $-10\,971$ \\
$-$ prix introuvable & $-14\,833$ & $-1\,106$ \\
$-$ déjà connue par ailleurs & --- & $-479$ \\
\midrule
\textbf{retenues} & \textbf{107\,976} & \textbf{13\,324} \\
membres concernés & 305 & 285 \\
\bottomrule
\end{tabular}
\end{center}
\vspace{1mm}
\pts{
\item Chaque ligne écartée a une raison affichée, et la somme est vérifiée automatiquement dans le notebook.
\item Les 479 transactions \og déjà connues \fg{} sont celles que le rapport annuel répète à l'identique.
\item Les deux sources ne se mélangent jamais sans le dire : on garde la trace de la provenance de chaque
transaction.
}
\src{source : cellules \og livre de comptes O/A/T \fg{} et \og collisions \fg{}. Le détail ligne à ligne des deux
cascades, avec la vérification de chaque somme, vit dans le notebook.}
\end{etapebox}}{\begin{etapebox}{8}{vert}{8 $\cdot$ ON RECONSTRUIT --- la photo devient un portefeuille}
\reponse{Une photo donne des valeurs en fourchettes, pas des quantités. On la convertit en parts détenues
au prix du jour de son dépôt, puis on applique les transactions.}
\phare{207}{portefeuilles de départ construits}{vert}
\vis{fig_10_entonnoir_h0.png}
\begin{center}\footnotesize
$h_0 = \dfrac{\tilde v}{P(\text{jour du dépôt})}$ \qquad achat : $h \mathrel{+}= \dfrac{\tilde v}{P}$
\qquad vente : $h \mathrel{-}= \min(q, h)$\\[0.5mm]
{\footnotesize $h$ = parts détenues, $P$ = prix du titre, $\tilde v$ = milieu de la fourchette déclarée,
$q$ = parts vendues}\\[0.8mm]
{\footnotesize\textbf{Sur un cas réel} : une position déclarée entre 1\,001 et 15\,000 \$, soit
8\,000 \$ au milieu, pour un titre à 59,60 \$ le jour du dépôt $\rightarrow$ \textbf{134 parts}.}
\end{center}
\vspace{0.8mm}
\pts{
\item On utilise le prix du jour du \textbf{dépôt}, jamais un prix antérieur : rien de ce qui suit ne
regarde le futur. 85,8\,\% de la valeur en actions de ces photos (543 M\$) devient ainsi des parts ; le
reste bute sur un prix manquant.
\item Injecter une photo n'ajoute ni ne retire de valeur : le jour de l'injection, le portefeuille ne
bouge pas d'un centime. C'est testé automatiquement.
\item Les membres sans photo commencent à zéro. Leur point de départ n'est pas nul, il est inconnu.
}
\src{source : cellules \og $h_0$ en parts \fg{} et \og livre de comptes photos \fg{}.}
\end{etapebox}}
\vspace*{\fill}
\newpage

\vspace*{\fill}
\rangee{\begin{etapebox}{9}{vert}{9 $\cdot$ ON VÉRIFIE --- les rapports annuels jugent}
\reponse{Le rapport annuel dit ce que l'élu possédait vraiment. On confronte notre reconstruction à cette
déclaration, position par position, à la date du dépôt.}
\phare{8 \% $\rightarrow$ 58 \%}{des positions déclarées retrouvées, sur la photo médiane}{vert}
\vis{fig_10_recouvrement.png}
\begin{center}\footnotesize
\begin{tabular}{@{}lccc@{}}
\toprule
sources utilisées & retrouvé & juste & valeur \\
\midrule
déclarations rapides seules & 8 \% & 59 \% & 54 \% \\
dossier complet & \textbf{58 \%} & 60 \% & 54 \% \\
\bottomrule
\end{tabular}\\[0.6mm]
{\footnotesize médianes par photo, sur 1\,376 et 1\,625 photos confrontables}
\end{center}
\vspace{1mm}
\pts{
\item \textbf{Valeur} : sur les positions communes, part de celles dont la valeur reconstruite tombe
dans la fourchette déclarée.
\item Le dossier complet ajoute la photo d'entrée et les transactions des rapports annuels. On retrouve
alors sept fois plus de positions sans perdre en justesse : 59 \% $\rightarrow$ 60 \%, et la valeur ne
bouge pas.
\item Une photo n'est confrontable que si une reconstruction existe à sa date : d'où 1\,376 et 1\,625
photos sur les 1\,945 exploitables. Même exercice sur les photos de sortie : 108 confrontées, 64 \%
retrouvé et 59 \% juste.
}
\src{source : cellules \og confrontation \fg{} et \og agrégats \fg{}.}
\end{etapebox}}{\begin{tcolorbox}[enhanced, breakable=false, colback=rouge!3, colframe=rouge!45, boxrule=0.5pt,
                  arc=2pt, left=2.2mm, right=2.2mm, top=1.8mm, bottom=1.8mm,
                  toptitle=0.7mm, bottomtitle=0.7mm, boxsep=0pt, before skip=1.9mm, after skip=0mm,
                  colbacktitle=rouge, coltitle=white, fonttitle=\footnotesize\bfseries,
                  title={CE QUE CETTE DONNÉE NE DIT PAS}]
\reponse{Six limites bornent tout ce qui précède. Elles ne sont pas des défauts de traitement : elles
tiennent à ce que la loi demande de déclarer, et à ce qu'elle n'exige pas.}
\pts{
\item Les valeurs sont déclarées en fourchettes, jamais au dollar près. Le milieu est une convention de
calcul, pas une mesure.
\item 181 membres arrivés n'ont déposé aucune photo d'entrée : leur point de départ reste inconnu.
\item Environ 42 \% des positions déclarées ne sont pas retrouvées, même avec le dossier complet. Elles
précèdent toute photo, ou portent sur des actifs sans ticker.
\item Les transactions issues des rapports annuels sont datées du jour de l'opération, mais publiées
373 jours plus tard : on les connaît bien après coup.
\item Le patrimoine sans ticker, soit 68\,266 M\$ déclarés, reste hors de portée du moteur.
\item Rien ici ne dit ce que valent ces trades : cette fiche décrit la donnée, et s'arrête là.
}
{\footnotesize\textbf{Pourquoi croire ces chiffres.} L'empreinte des tables est conforme au manifeste.
Les mandats sont re-croisés avec le référentiel officiel (98,9 \%). Chaque cascade et chaque tableau
croisé sont vérifiés automatiquement. La classe d'un actif est décidée avant tout prix. La reconstruction
est testée de trois façons indépendantes. Aucun chiffre n'est écrit en dur : tout se régénère en
ré-exécutant le notebook.}
\end{tcolorbox}}
\vspace*{\fill}

\end{document}
